% !TeX root = ../tg.tex

\chapter{课题研究的背景和内容}
\section{研究背景}
随着边缘计算与隐私保护需求的日益增长，分布式机器学习框架正逐步成为解决大规模数据协同训练的重要手段。近年来，联邦学习（Federated Learning, FL）因其能够在不共享原始数据的前提下实现多方协作建模而受到广泛关注 \cite{b1}。然而，传统联邦学习要求每个客户端独立训练整个模型，这对资源受限的边缘设备而言是一项沉重的负担。

在此背景下，分割-联邦学习（SplitFed）应运而生 \cite{thapa2022splitfed}。该方法结合了分割学习（Split Learning, SL） \cite{b2} 的思想，通过在客户端与服务器之间对神经网络模型进行纵向拆分，实现了更高效的分布式训练。具体而言，在分割-联邦学习中，客户端仅负责前向传播和部分反向传播的浅层模型计算，而深层模型则由服务器统一维护与更新。这种设计不仅减轻了客户端的计算压力，还有效保护了用户隐私——因为数据特征始终保留在本地，不会被上传至服务器或其他客户端。

尽管分割-联邦学习在降低通信开销和提升系统效率方面展现出显著优势，其在实际部署过程中仍面临诸多挑战。首先，分割-联邦学习依赖于各客户端提供的训练数据质量。理想情况下，这些数据应具有准确无误的标签，并且在不同客户端间呈现独立同分布（IID）特性。然而，在真实应用场景中，由于缺乏集中式的数据预处理机制，客户端数据往往包含大量错误标签（即标签噪声），同时其标签分布也呈现出高度的非独立同分布（Non-IID）特性。这类问题不仅会严重削弱模型的泛化能力，还会拖慢模型收敛速度，甚至导致训练过程陷入局部最优。

其次，分割-联邦学习虽然将主要的训练任务交由服务器承担，但客户端仍需参与每一轮训练中的激活值传递与梯度计算。对于算力有限、内存紧张、带宽受限的边缘设备而言，这依然是不小的负担。尤其是在面对复杂模型结构时，客户端加载完整的计算图可能成为系统性能的瓶颈，限制了分割-联邦学习在低功耗设备上的广泛应用。

基于上述分析，本课题旨在围绕分割-联邦学习框架中因数据质量问题与硬件资源约束所引发的核心挑战，开展系统性的研究工作。具体目标如下：

\begin{enumerate}
    \item \textbf{深入研究异构数据对分割-联邦学习的影响以及优化方案}：通过构建数学模型，揭示标签噪声与非独立同分布数据对模型训练动态的影响机制，为后续算法优化提供理论依据。
    同时在保持分割-联邦学习轻量化特性的基础上，探索增强模型在存在噪声与非均衡数据条件下的稳定性和泛化性能的途径；
    % \item \textbf{设计鲁棒性强、适应性广的SplitFed训练算法}：
    \item \textbf{探索分割-联邦学习中轻量化客户端训练方案}：引入无反向传播的梯度估计方法，以显著降低客户端侧的计算复杂度、内存占用与通信成本，从而提升整个系统的可扩展性与部署灵活性。
\end{enumerate}

通过以上研究，期望能够推动分割-联邦学习在资源受限与数据不确定场景下的实用化进程，为其在医疗、金融、智能物联网等高隐私敏感领域的广泛应用奠定坚实基础。
% \section{课题研究的基础}

\section{研究内容}
本研究围绕分割-联邦学习展开，聚焦于两个核心问题：异构数据对分割-联邦学习性能的影响及优化方案 ，以及轻量化的客户端训练机制设计与实现 。具体研究内容包括以下三个方面：
\subsection{收敛性分析与算法设计：异构数据对分割-联邦学习的影响以及优化方案}
\subsubsection{分割-联邦学习的收敛性分析}
在分割-联邦学习框架下，模型被纵向划分并分别部署于客户端与服务器端，客户端仅负责浅层网络的计算，而深层网络由服务器统一维护与更新。这种架构虽降低了客户端的计算负担，但也引入了新的挑战，尤其是在面对非独立同分布（Non-IID）数据和标签噪声时，其收敛行为尚不明确。

本研究系统地分析了分割-联邦学习在均方误差（MSE）和交叉熵损失函数下的收敛特性。通过建立理论框架，推导了在存在标签噪声的情况下，客户端本地随机梯度的方差上界及其与期望梯度之间的偏差。在此基础上，进一步探讨了不同数据分布、模型拆分点、本地训练轮次等因素对全局模型收敛速度和最终性能的影响。实验验证表明，在合理配置训练参数的前提下，分割-联邦学习可以在标签噪声和非独立同分布数据下稳定收敛，并保持较高的模型精度。

该部分的研究为后续鲁棒性算法的设计提供了理论支撑，也为分割-联邦学习在资源受限设备上的应用奠定了基础。
\subsubsection{对标签噪声和非独立同分布数据具有鲁棒性的分割-联邦学习算法}
由于实际应用场景中数据往往具有异构性和噪声特征，如何提升分割-联邦学习对此类数据的鲁棒性成为关键问题。为此，本研究提出了 SFL-DCT 算法，旨在有效识别并修正标签噪声，同时缓解非独立同分布数据带来的负面影响。

SFL-DCT 分为三个阶段：第一阶段通过局部迭代训练提取每个客户端的数据特征，并基于 LID（Local Intrinsic Dimensionality）或 LE（Label Entropy）指标检测可能含有噪声标签的客户端；第二阶段利用干净客户端集合训练一个鲁棒模型，并用于识别噪声样本并进行标签修正；第三阶段则在修正后的数据集上进行标准的分割-联邦学习训练。整个过程结合高斯混合模型（GMM）分类技术，实现了对噪声客户端和噪声样本的自动识别与处理。

实验结果表明，SFL-DCT 在 CIFAR-10、CIFAR-100 和 Fashion-MNIST 等多个数据集上均表现出良好的鲁棒性。即使在高达 60\% 的标签噪声率下，SFL-DCT 仍能维持较高的测试准确率，并优于传统联邦学习中的噪声鲁棒算法（如 FedCorr 和 FedNoRo）。此外，在 Non-IID 场景下，SFL-DCT 显著提升了模型的泛化能力，证明了其在异构环境中的有效性。
\subsection{轻量化客户端训练方案：无反向传播训练在分割-联邦学习中的应用}
为了进一步降低客户端的计算和内存开销，本研究首次将无反向传播训练机制引入分割-联邦学习框架中。传统的分割-联邦学习客户端需要执行前向传播和部分反向传播操作，这在资源受限设备上仍具有一定负担。为此，我们探索了多种无需依赖完整梯度反传的训练方法，如直接反馈对齐（Direct Feedback Alignment）、固定反馈权重等策略，并将其适配到分割-联邦学习的训练流程中。

具体而言，客户端侧采用无反向传播机制完成局部模型更新，仅需前向计算与少量反馈信号即可调整参数。这种设计显著减少了内存占用和计算复杂度，同时降低了通信成本。实验结果显示，相较于基于标准 SGD 的分割-联邦学习方法，引入无反向传播机制后，客户端的训练时间平均减少约 30\%，且在多个视觉任务上的模型精度下降幅度控制在 2\% 以内，展现了良好的性能与效率平衡。

该部分工作不仅拓展了分割-联邦学习的适用场景，还为边缘设备上的低功耗训练提供了新思路，具有较强的工程实用价值。